% Capítulo "Considerações finais" do relatório -- EDITÁVEL À MÃO (specs/2026-11-05_versao_final_site R5).
% Exportado uma vez por gera_latex.py --exportar-capitulos (2026-11-05); desde então o TEXTO é escrito aqui, não
% vem mais de relatorio/textos_curados.json (que é só do site). Os DADOS continuam gerados: as figuras
% (\figura{...}{../../visualizacoes/a4/...} e mapas/a4/) saem do analise.py; as tabelas no corpo
% (\input{gerado/tabelas_corpo/...}) e os rótulos do apêndice (\ref{tab:...}) saem do gera_latex.py, que confere
% a cada execução se tudo o que este arquivo cita existe.

\chapter{Considerações finais}\label{cap:consideracoes}

A análise dos indicadores aqui trabalhados nos dá uma visão multidimensional da primeira infância na cidade do Rio de Janeiro. Ela mostra o que melhorou, mas também as desigualdades e os desafios que a prefeitura precisa acompanhar de perto. Fica claro que a situação das crianças não se resume a um único fator: saúde, renda, educação, alimentação, vacinas e o território onde moram estão todos interligados.

Na saúde, vimos uma queda importante nas mortes infantis por causas evitáveis. Porém, muitas dessas mortes ainda estão ligadas à falta de cuidado adequado com a mulher durante a gravidez. Como a maioria acontece nos primeiros meses de vida do bebê, é fundamental melhorar a rede de apoio desde o pré-natal até o primeiro ano da criança, incluindo o momento do parto.

Além disso, a realidade muda muito dependendo da região da cidade. As Zonas Norte e Oeste concentram mais famílias no Cadastro Único (CadÚnico), mais bebês nascendo com baixo peso e mais casos de mortalidade. Isso aponta as áreas onde a prefeitura precisa investigar mais a fundo as condições sociais e reforçar os serviços oferecidos à população.

O CadÚnico reforça o peso da questão financeira. A grande concentração de famílias com crianças pequenas nas faixas de renda mais baixas mostra que o planejamento público deve considerar a multidimensionalidade da pobreza. É preciso unir ações de saúde, educação e assistência social, focando especialmente nessas áreas mais vulneráveis.

Na educação, as crianças passam a ir mais à escola à medida que crescem, mas ainda há quem esteja fora da pré-escola, que tem a obrigatoriedade com força constitucional a partir dos 4 anos. O desafio aqui não é só olhar para a média da cidade, mas descobrir em quais bairros essas crianças estão e quais obstáculos as impedem de estudar.

A vacinação é outro ponto de alerta. Depois de uma queda entre 2019 e 2022, as taxas começaram a subir em 2023. O problema é que essa melhora não foi igual para todas as vacinas, e ainda há atrasos nas doses de reforço. Já os dados de nutrição revelam duas faces do mesmo problema: a desnutrição persiste, ao mesmo tempo em que o sobrepeso e a obesidade infantil vêm crescendo. Isso exige um olhar mais amplo sobre a alimentação, que não trate o problema apenas como falta de comida ou apenas como excesso de peso, a segurança alimentar também considera a nutrição e qualidade dessa alimentação.

Tudo isso mostra que os desafios da primeira infância no Rio são, simultaneamente, de saúde, educação, nutrição e moradia. Para dar conta, o poder público não deve atuar de forma isolada, e sim integrar o trabalho de suas diferentes secretarias e serviços.

Ainda assim, é preciso olhar para os números com cautela. As informações vêm de bases de dados diferentes e não indicam necessária causalidade.  O grande valor desse diagnóstico é colocar uma lupa de forma integrada nas tendências e desigualdades do território no que diz respeito à primeira infância. Esse levantamento deve ajudar a prefeitura a planejar, agir e avaliar constantemente suas ações.
